\section{Bounded testing experiment and geometry proofs}

\subsection{The bounded testing construction}

The minimax converse in \cref{obj:thm:adaptive-minimax} is supported by a pair of causal laws with two-valued potential outcomes. The construction places a smooth response perturbation around the evaluation point and uses a common radial propensity to control the treatment information available there. Its radial distribution accounts for intersection with the covariate cube, so the experiment accommodates evaluation points on faces and corners as well as in the interior. The following definition specifies both the observed laws and their causal completions.

\begin{definitionv}{def:bounded-lower-pair}[Bounded radial testing pair]
The bounded testing construction is the measure-valued family \(P_{j,h}\) on observations \((X,A,Y)\), indexed by \(j\in\{0,1\}\) and \(n\in\mathbb N\), with parameters \(d\in\mathbb N\), \(\beta,B,L,\gamma\in\mathbb R\), \(x_0\in\mathbb R^d\), and constants \(a,c>0\). Set
\[
\mathcal X:=[0,1]^d,\qquad
h:=c h_{\gamma,n},\qquad
\theta_0:=\min\{1/4,L/(4B)\},
\]
where \(h_{\gamma,n}\) is the oracle mesh for the smoothness and overlap parameters. Define the radial distance, cube-truncated radial function, and common propensity by
\[
R(x):=\|x-x_0\|_\infty,\qquad
F_R(t):=\prod_{i=1}^d
\max\{0,\min\{1,x_{0i}+t\}-\max\{0,x_{0i}-t\}\},
\qquad
e_0(x):=F_R(R(x))^{1/(\gamma-1)}.
\]
For \(x_0\in\mathcal X\) and \(t\ge0\), \(F_R(t)\) is the probability of \(R(X)\le t\) under uniform covariates. Define the smooth product bump and treated success parameters by
\[
\psi(z):=\prod_{i=1}^d
\begin{cases}
\exp\{1-(1-z_i^2)^{-1}\},& |z_i|<1,\\
0,& |z_i|\ge1,
\end{cases}
\qquad
p_{0,h}(x):=\theta_0,\qquad
p_{1,h}(x):=\theta_0+ah^\beta\psi((x-x_0)/h).
\]

The totalized construction uses the positive-part weights
\[
w_0(p):=\max\{1-p,0\},\qquad
w_1(p):=\max\{p,0\}.
\]
For each covariate \(x\), take the product of the three two-point measures with weights
\[
\begin{array}{c|cc}
 & \text{value }0 & \text{value }1\\ \hline
A & w_0(e_0(x)) & w_1(e_0(x))\\
Y(0)/B & w_0(\theta_0) & w_1(\theta_0)\\
Y(1)/B & w_0(p_{j,h}(x)) & w_1(p_{j,h}(x))
\end{array}
\]
where the potential-outcome measures place their masses at \(0\) and \(B\), respectively. Integrate this product measure against Lebesgue measure restricted to \(\mathcal X\), and let \(P_{j,h}\) be its image under
\[
(X,A,Y(0),Y(1))
\longmapsto
\bigl(X,A,AY(1)+(1-A)Y(0)\bigr).
\]

The admissible domain consists exactly of the constants \(a,c>0\) for which
\[
e_0(x),p_{0,h}(x),p_{1,h}(x)\in[0,1]
\]
for every \(n\ge1\) and every \(x\in\mathcal X\), with \(h=c h_{\gamma,n}\). On this domain, \(P_{j,h}\) is the observed law obtained from uniform \(X\) and conditionally mutually independent \(A,Y(0),Y(1)\) given \(X\), with
\[
A\mid X\sim\operatorname{Bernoulli}(e_0(X)),\qquad
Y(0)\mid X\sim B\operatorname{Bernoulli}(\theta_0),\qquad
Y(1)\mid X\sim B\operatorname{Bernoulli}(p_{j,h}(X)),
\qquad
Y=AY(1)+(1-A)Y(0).
\]
For fixed admissible \(a,c\), define the sample-size-indexed collection
\[
\mathcal P_{\mathrm{pair}}
:=\{(P_{0,h},P_{1,h}):n\ge1,\ h=c h_{\gamma,n}\}.
\]
\end{definitionv}

The positive-part weights in \cref{obj:def:bounded-lower-pair} specify a common measure-family construction across parameter values. On the admissible domain, each conditional two-point measure has total mass one and the construction yields the displayed Bernoulli probability laws. The sample-size-indexed collection \leanref{sym:Ppair}{\ensuremath{\mathcal P_{\mathrm{pair}}}} uses fixed amplitude and mesh constants; sample size changes the perturbation through the oracle mesh.

Admissibility supplies the probability interpretation, while membership in the model of \cref{obj:def:model} additionally requires the response and overlap restrictions. The next result makes these requirements compatible with small statistical information and separation of the chosen response representatives at the evaluation point.

\begin{lemmav}{lem:bounded-lower-pair}[Bounded testing pair guarantees]
Suppose that:
\begin{itemize}
\item \textbf{(Dimension and exponents.)} \(d\in\mathbb N\), \(d>0\), \(\beta>1\), and \(\gamma>1\).
\item \textbf{(Response constants.)} \(B>0\) and \(L>0\).
\item \textbf{(Model constants.)} \(C\ge1\) and \(0<c_f\le1\).
\item \textbf{(Evaluation point.)} \(x_0\in[0,1]^d\).
\end{itemize}
There exist constants \(a,c,c_0>0\), fixed across sample sizes, such that, for every integer \(n\ge1\), with \(h=ch_{\gamma,n}\), where \(h_{\gamma,n}\) is the oracle mesh, the construction in \cref{obj:def:bounded-lower-pair} satisfies
\[
e_0(x),\ p_{0,h}(x),\ p_{1,h}(x)\in[0,1]
\qquad\text{for every }x\in[0,1]^d.
\]
Both observed laws \(P_{0,h}\) and \(P_{1,h}\) belong to \(\mathcal P_\gamma\), the model class with parameters \(d,\beta,B,L,C,c_f,\gamma\). Under each of their constructed causal completions,
\[
Y(0),Y(1)\in\{0,B\}
\qquad\text{almost surely}.
\]
Moreover, \(P_{0,h}\ll P_{1,h}\), the log likelihood ratio
\(\log(dP_{0,h}/dP_{1,h})\) is integrable under \(P_{0,h}\), and
\[
n\,\operatorname{KL}(P_{0,h},P_{1,h})\le\frac18.
\]
For the chosen Hölder response representatives \(\mu_{1,P_{0,h}}\) and \(\mu_{1,P_{1,h}}\) associated with their model-class memberships,
\[
\left|\mu_{1,P_{1,h}}(x_0)-\mu_{1,P_{0,h}}(x_0)\right|
\ge c_0 r_{\gamma,n},
\qquad
r_{\gamma,n}=h_{\gamma,n}^{\beta},
\]
where \(r_{\gamma,n}\) is the oracle response-recovery risk scale.
\end{lemmav}

\Cref{obj:lem:bounded-lower-pair} supplies the three components of the testing comparison: two laws within the fixed-constant model, a bound on their relative entropy, and response separation at the oracle scale. Here \(\operatorname{KL}\) denotes Kullback--Leibler relative entropy. The absolute-continuity and integrability clauses support the information comparison for independent samples, whose laws are \(P_{0,h}^{\otimes n}\) and \(P_{1,h}^{\otimes n}\). The chosen Hölder representatives specify the pointwise targets, including their values on the cube boundary.

The lower guarantee in \cref{obj:thm:adaptive-minimax} uses this pair through the two-point testing principle \citep[Section~2.2]{tsybakov2009}. The information bound and target separation provide the comparison for expected absolute error within the subclass admitting potential outcomes in \(\{0,B\}\). The risk ordering in \cref{obj:thm:adaptive-minimax} then places the same oracle scale below the pointwise and expected spatial-supremum minimax risks of \cref{obj:def:risks}.

The quantifiers of the construction preserve the pointwise scope of the converse. For each fixed \(\gamma>1\) and each fixed \(x_0\) in the closed cube, the constants are chosen once and apply to every integer \(n\ge1\). Their permitted dependence on the evaluation point and fixed model parameters matches \cref{obj:thm:adaptive-minimax}. Cube truncation in \cref{obj:def:bounded-lower-pair} and the intrinsic smoothness convention in \cref{obj:def:holder} keep boundary evaluation within this same comparison.

\subsection{The thin-strip geometry}

The geometry assertions in \cref{obj:prop:strict-enlargement} concern the strip law of \cref{obj:def:thin-strip}. That construction uses the cube center and a radial-plus-strip propensity; its strip exponent is the parameter specified there, separately from the amplitude constant in the testing pair. Constant potential-outcome means keep response smoothness fixed while the propensity determines the local treated design.

The model-membership and tail clauses of \cref{obj:prop:strict-enlargement} place this law within the global-tail causal class. Its remaining clauses describe two distinct aspects of local geometry. The anti-concentration statement concerns the displayed selected propensity and the pointwise supremum over Euclidean neighborhoods, as specified in \cref{obj:def:dorn-a3}. The numerical clause has a counterexample even when propensity values equal to one are permitted. This establishes the selected-version failure underlying the strict enlargement within the causal class.

The second-moment clause instead concerns treatment-weighted variation in supremum-norm neighborhoods of the cube center. Its denominator is the probability of both treatment and neighborhood membership, so the ratio measures the conditional second moment of the second coordinate after normalization by neighborhood width. The zero limit in \cref{obj:prop:strict-enlargement} records degeneration of that coordinate's local treated variation under the stated strip-exponent range.

These conclusions have complementary roles. The selected-version statement establishes the comparison with the local anti-concentration predicate, while the normalized moment describes the shape of the treated distribution itself. Together with model membership and the global tail inequality, they characterize the treatment geometry covered by \cref{obj:thm:adaptive-minimax}.
